&\hspace{0.43cm}\;-\; \sum_i \big(A_{i,zz} I_{iz}+A_{i,zx} I_{ix}+A_{i,zy} I_{iy}\big) \label{eq:H-1_def}
\end{aligned}
\end{equation}
By contrast, in the $\ket{0_e}$-manifold the projected $S_z$ contribution vanishes, since $S_z\ket{0_e}=0$, so that the nuclei evolve only under their Zeeman $\gamma_i B_0 I_{iz}$ and quadrupolar terms $Q_i\!\left(I_{iz}^2-\tfrac{2}{3}\right)$,
\begin{equation}
\begin{aligned}
H_{0} &= -\sum_i \gamma_i B_0 I_{iz} \;+\; \sum_i Q_i\!\left(I_{iz}^2-\tfrac{2}{3}\right)
\label{eq:H0_def}
\end{aligned}
\end{equation}
Recombining the secular electronic Hamiltonian with the two conditional nuclear Hamiltonians yields the block form
\begin{equation}
H_{\mathrm{NV}}
= H_e
+ \ket{0_e}\!\bra{0_e}\otimes H_{0}
+ \ket{-1_e}\!\bra{-1_e}\otimes H_{-1}.
\label{eq:block_form}
\end{equation}

\paragraph{Diagonalizing the Nuclear Basis.}
We want to diagonalize the $H_{-1}$ Hamiltonian because it contains off-diagonal terms with respect to the electronic quantization axis, in contrast to the $H_{0}$ Hamiltonian.
In the $H_{-1}$ Hamiltonian, the transverse hyperfine components tilt the effective nuclear field away from the NV axis. For nucleus $i$, the combined secular hyperfine and nuclear Zeeman contributions in the $H_{-1}$ Hamiltonian can be grouped as
\begin{equation}
    H^{(i)}_{-1,\mathrm{hf}} = (A_{i,zz}+\gamma_iB_0)I_{iz}
    +A_{i,zx}I_{ix}
    ++A_{i,zy}I_{iy}
    \label{eq:Hminus_hf}
\end{equation}
which, by itself, can be geometrically interpreted as an interaction of the nuclear spin with an effective static field
\[
\vec b_i = 
\begin{pmatrix}
A_{i,zx}\\
A_{i,zy}\\
A_{i,zz}+\gamma_iB_0\\
\end{pmatrix}
\]
whose orientation corresponds to the current quantization axis in the $H_{-1}$ Hamiltonian.


In order to ultimately rotate the basis of $H_{-1}$, we parameterize the orientation of $\vec b_i$ by spherical angles $\theta_i,\phi_i$ \cite{schweiger2001principles},
\begin{equation}
\tan\phi_i = \frac{A_{i,zy}}{A_{i,zx}},
\qquad
\tan\theta_i = \frac{\sqrt{A_{i,zx}^2+A_{i,zy}^2}}
{A_{i,zz}+\gamma_iB_0},
\label{eq:angles}
\end{equation}
and define its magnitude
\begin{equation}
\omega_i = |\vec b_i| =
\sqrt{ (A_{i,zz}+\gamma_iB_0)^2
+ A_{i,zx}^2
+ A_{i,zy}^2 }.
\label{eq:omega_i}
\end{equation}

Based on these spherical angles, we introduce a rotated nuclear basis $(x'_i,y'_i,z'_i)$ whose quantization axis $z'_i$ is aligned with $\vec b_i$. The $I_{iz}$ operator in this rotated basis basis becomes 
\begin{equation}
    I'_{iz} = I_{ix}\sin\theta_i\cos\phi_i 
    + I_{iy}\sin\theta_i\sin\phi_i 
    + I_{iz}\cos\theta_i.
    \label{eq:Izprime}
\end{equation}
With this choice, the transverse contribution in the $H_{-1}$ Hamiltonian become diagonal:
\[
H^{(i)}_{-1,\mathrm{hf}}= \vec b_i\cdot\vec I_i = \omega_i I'_{iz}.
\]

The $H_{-1}$ Hamiltonian therefore takes the form
\begin{align}
H_{-1}
&=
-\sum_i\omega_i I'_{iz}
+
\sum_iQ_i\!\left(I_{iz}'^{\,2}-\tfrac23\right).
\label{eq:H-1_rot}
\end{align}

Thus, after rotating the nuclear basis, the quantization axes associated with the $H_{-1}$ Hamiltonian are collected onto the effective hyperfine axis of each nucleus, rendering the Hamiltonian diagonal in the rotated basis.\\

\paragraph{Electronic transitions between nuclear quantization frames.}

While the $H_0$ Hamiltonian remains quantized along the NV axis, the $H_{-1}$ Hamiltonian is now quantized along the effective hyperfine axis defined by $\vec b_i$. The two nuclear Hamiltonians therefore have a relative tilt between their quantization frames.

The rotated operator \(I'_{iz}\) implicitly defines a nuclear rotation
\(R_i\) satisfying
\[
I'_{iz}=R_i I_{iz} R_i^\dagger.
\]
Since the effective field \(\vec b_i\) is parameterized by the spherical angles \((\theta_i,\phi_i)\), this rotation is generated by
\[
R_i
=
\exp\!\left[
-i\theta_i
\left(
-I_{ix}\sin\phi_i
+
I_{iy}\cos\phi_i
\right)
\right],
\]
with collective rotation
\[
R=\prod_iR_i.
\]

For the weakly misaligned nuclei considered here,  \(\theta_i\ll1\), the rotation may be expanded to first order,
\[
R
\simeq
1-iG,
\]
where
\[
G
=
\sum_i\theta_i
\left(
-I_{ix}\sin\phi_i
+
I_{iy}\cos\phi_i
\right).
\]

Rather than rewriting both nuclear Hamiltonians in a common operator basis, we absorb this relative nuclear rotation into the electronic transition operators. Defining the electronic raising and lowering operators as
\[
\sigma_+
=
\ket{0_e}\!\bra{-1_e},
\qquad
\sigma_-
=
\ket{-1_e}\!\bra{0_e},
\]
the microwave interaction
\[
H_{\mathrm{MW}}
=
\frac{\gamma_eB_{1,\mathrm{MW}}(t)}{\sqrt2}
(\sigma_+ + \sigma_-)
\]
therefore becomes
\begin{equation}
\begin{aligned}
H_{\mathrm{MW}}
&=
\frac{\gamma_eB_{1,\mathrm{MW}}(t)}{\sqrt2}
\left(
\sigma_+\otimes R^\dagger
+
\sigma_-\otimes R
\right)
\\
&\approx
\frac{\gamma_eB_{1,\mathrm{MW}}(t)}{\sqrt2}
\left[
\sigma_x
-
\sigma_y G
\right].
\end{aligned}
\label{eq:HMW_rot}
\end{equation}

The first term describes the bare electronic microwave transition, while the second term is a first-order electron-nuclear modulation arising from the mismatch between the nuclear quantization frames of the two electronic manifolds.\\




\paragraph{Interaction picture.}

The next step of the derivation is a transformation to the interaction picture that removes the free precession dynamics and leaves only the control-induced dynamics explicit. We therefore decompose the Hamiltonian as
\begin{equation}
H_{\mathrm{NV}}
=
H_{\mathrm{free}}
+
H_{\mathrm{MW}},
\label{eq:H_split}
\end{equation}
where the free evolution is generated by the diagonal electronic and nuclear
terms,
\begin{equation}
\begin{aligned}
H_{\mathrm{free}}
&=
\frac{\Lambda_s}{2}\sigma_z
\\
&\quad
+
\ket{0_e}\!\bra{0_e}\otimes
\left[
-\sum_i\gamma_iB_0 I_{iz}
+
\sum_iQ_i\!\left(I_{iz}^{\,2}-\tfrac23\right)
\right]
\\
&\quad
+
\ket{-1_e}\!\bra{-1_e}\otimes
\left[
-\sum_i\omega_i I'_{iz}
+
\sum_iQ_i\!\left(I_{iz}'^{\,2}-\tfrac23\right)
\right],
\end{aligned}
\label{eq:H_free}
\end{equation}
while the microwave interaction is
\begin{equation}
\begin{aligned}
H_{\mathrm{MW}}
&\approx
\frac{\gamma_e B_{1,\mathrm{MW}}(t)}{\sqrt2}
\left[
\sigma_x
-
\sigma_y G
\right]\,.
\end{aligned}
\label{eq:HMW_rot}
\end{equation}


We choose \(H_{\mathrm{free}}\) as the generator of the interaction picture, such that
\begin{equation}
H_{\mathrm{NV}}^{(I)}(t)
=
e^{iH_{\mathrm{free}}t}
H_{\mathrm{MW}}
e^{-iH_{\mathrm{free}}t}.
\label{eq:IP_transform}
\end{equation}

We first transform the bare electronic part of the microwave interaction,
\[
H_{\mathrm{MW},x}
=
\frac{\gamma_e B_{1,\mathrm{MW}}(t)}{\sqrt2}\sigma_x .
\]
Its interaction-picture form is
\begin{equation}
H_{\mathrm{MW},x}^{(I)}(t)
=
\frac{\gamma_e B_{1,\mathrm{MW}}(t)}{\sqrt2}
e^{iH_{\mathrm{free}}t}
\sigma_x
e^{-iH_{\mathrm{free}}t}.
\label{eq:HMWx_IP_start}
\end{equation}


Since \(H_{\mathrm{free}}\) is diagonal in the electronic basis and in the natural nuclear basis of each manifold, its propagator acts only through energy-dependent phase accumulation,
\[
e^{iH_{\mathrm{free}}t}
=
\sum_{e,\mathbf m}
e^{iE_e(\mathbf m)t}
\ket{e,\mathbf m}\!\bra{e,\mathbf m}.
\]
Using \(\sigma_x=\sigma_++\sigma_-\), the interaction picture evolution of the electronic transition operator is therefore governed by the energy difference between the connected free eigenstates,
\begin{equation}
\begin{aligned}
&e^{iH_{\mathrm{free}}t}
\sigma_x
e^{-iH_{\mathrm{free}}t}
\notag\\
&=
\sum_{\mathbf m}
\Big[
\sigma_+
e^{i[E_0(\mathbf m)-E_{-1}(\mathbf m)+\Lambda_s]t}
\\
&\hspace{1.5cm}
+
\sigma_-
e^{-i[E_0(\mathbf m)-E_{-1}(\mathbf m)+\Lambda_s]t}
\Big]
\otimes
\ket{\mathbf m}\!\bra{\mathbf m}.
\end{aligned}
\label{eq:sigmax_IP_phase}
\end{equation}
Defining the configuration-dependent transition frequency
\begin{equation}
\begin{aligned}
\Lambda(\mathbf m)
&:=
\Lambda_s
+
E_0(\mathbf m)
-
E_{-1}(\mathbf m)
\\
&=
\Lambda_s
+
\left[
